% Ad Atticum 13.40 --- headnote
% ad-atticum-13-40, 8 August 45 BC, Tusculum

\textit{Cicero to Atticus, written at the Tusculan villa on 7 or 8
August 45 BC --- Perseus dateline \textit{Scr. in Tusculano vii
aut vi Id. Sext. a. 709 (45)}. The most textually corrupt of
this Tusculan cluster: two short sections studded with daggered
cruxes that the manuscript tradition cannot disentangle. The
nephew is the subject again. Brutus has reported to Cicero that
young Quintus is now declaring himself for the \textit{boni} ---
the loyalists, the old republican set --- and Cicero's response
is acid: \textit{good news indeed} (in Greek), but where will he
find any? Unless he has hanged himself in the meantime. A second
Greek word, \textit{philotechnema}, recalls a tableau Cicero saw
in the Parthenon during his youth --- Ahala and Brutus, the
tyrannicides of legend --- as the appropriate emblem of what the
younger Quintus is now pretending to.}

\textit{The second section is the practical one: should Cicero
come to Rome or stay put? He confesses he has been
\textit{taken in} (\textit{kekepphomai} --- a colloquial
Greek verb borrowed from the gull, the bird that is easy to
catch). He wants Atticus's reading of \textit{the whole
picture} (another Greek tag) at first light tomorrow. The
register is hurried and elliptical; the daggered cruxes are
kept in place, and the translation does not invent through
them.}
